\section{Theoretical and Computational Framework}
\subsection{Representational levels}
The model represents each entry in a unified lexicon at three complementary levels: the conceptual level, the lemma level, and the phonological level. For each lexical entry $i$, a conceptual vector $\mathbf c_i$, a lemma vector $\boldsymbol\ell_i$, and a phonological form vector $\mathbf p_i$ represent the corresponding concept, lemma, and phonological form of the same word. All vectors have dimensionality $D$. Each level uses independently generated base vectors to construct its representations within a level-specific coordinate system. The correspondence between the three representations of a lexical entry is specified through explicit cross-level associations. In the current implementation, the lemma level serves as a lexical identity interface; shared grammatical features, inflectional structure, and syntactic processing remain topics for future extensions.

Representations at the conceptual and phonological levels are structured through relations to other representations within the same level. For example, the phonological form of a word is related to representations of its phonological components; conceptual representations similarly incorporate associations with other conceptual representations. These relations are specified as graphs: nodes identify the represented concepts, phonological forms, or their components, and edges specify the relations between them. Each node is assigned a vector that incorporates its relations to other nodes, yielding a stored vector description of the lexical structure. These vectors are used to construct the neural network. During simulation, populations of spiking neurons encode time-dependent approximations to the vectors in their collective activity. The model thus connects two descriptions: a specification of lexical structure in terms of graphs and vectors, and its neural realization in terms of spiking population activity. Graph edges contribute to the construction of the stored vectors, while synaptic connections implement the transformations between neurally represented states.

\subsection{Neural representation and transformation}
The NEF describes how populations of neurons represent quantities and approximate transformations of those quantities \citep{eliasmith2003}. A useful rate-level description of the encoding of a vector $\mathbf x$ is
\begin{equation}
 a_j(\mathbf x)=G_j\!\left(\alpha_j\mathbf e_j^{\mathsf T}\mathbf x+J_j^{\mathrm{bias}}\right),
\end{equation}
where $\mathbf e_j$ is an encoder, $\alpha_j$ a gain, $J_j^{\mathrm{bias}}$ a bias current, and $G_j$ the neuronal response function. The simulations use spiking leaky integrate-and-fire (LIF) neurons, rather than substituting these rates for spike generation. A decoded state can be written
\begin{equation}
 \hat{\mathbf x}(t)=\sum_j\mathbf d_j\,(h*s_j)(t),
 \label{eq:decode}
\end{equation}
where $s_j$ is a spike train, $h$ a filtering kernel, and $\mathbf d_j$ a decoder. Decoders are computed to approximate the specified represented quantity or transformation over sampled evaluation points. This construction is distinct from learning lexical associations through experience.

The model is implemented in Nengo \citep{bekolay2014} with NengoSPA components. State vectors are distributed over one-dimensional ensembles, so Eq.~\eqref{eq:decode} applies componentwise within the corresponding ensemble array. The architecture specifies a functional computation; it does not provide a one-to-one anatomical assignment for each array.

\subsection{Semantic pointers and structured representations}
The SPA uses high-dimensional vectors to support both distributed representation and symbol-like operations \citep{eliasmith2013}. In the present model, three operations construct a structured vector from a node's base vector and its relations to other nodes: circular convolution, addition, and normalization. Circular convolution binds a vector specifying a relation to the vector representing its target. For example, binding the relation vector for \texttt{AssocWith} to a handle-related conceptual vector encodes the contribution ``associated with a handle.'' The resulting vector has the same dimensionality $D$ as its inputs and expresses the target in the context of that relation. Binding thus produces a new pattern within the same vector space, with a direction typically distinct from those of either input.

Addition combines these bound contributions with the node's own base vector. For example, the conceptual representation of \emph{Tasse} combines its base vector with contributions encoding associations with a dish-related concept and a handle-related concept. This sum forms a single vector containing contributions from several aspects of the word's conceptual structure. Normalization then rescales the sum to length 1 while preserving its direction. The stored representations therefore lie on the unit sphere in the $D$-dimensional space. A successfully retrieved neural state approximates one of these unit-length vectors; its decoded length can vary during activation buildup and through fluctuations in spiking activity.

Structured knowledge representation provides a broader context for combining relations and their targets in vector form \citep{crawford2016}. A central consequence in the present model is the emergence of representational similarity from shared structure. For example, \emph{Tasse} and \emph{Kanne} both include the same handle-related bound contribution. Such shared contributions can bring their normalized vectors closer together on the unit sphere. Groups of representations with overlapping structure can consequently occupy nearby regions of this sphere, with their precise positions determined by the full combination of base vectors and bound contributions. In this way, the graph structure shapes the geometry of conceptual and phonological similarity within each level.

In the present model, the lexical vectors are specified computationally and then realized in populations of spiking neurons. Their construction starts from a predefined lexical graph and randomly generated base and relation vectors, which are combined using the encoding operations described above. The resulting representations therefore reflect both the supplied lexical structure and the chosen encoding scheme. This dependence on the chosen lexical specification is a general property of the modeling approach and applies equally to alternative lexicons. Establishing how closely the resulting similarities correspond to human semantic or phonological similarity judgments requires an independent empirical comparison. This introduces a degree of arbitrariness in the assignment of vectors and constitutes a limitation of the current model: lexical representations and their associations are specified in advance rather than acquired through learning from linguistic experience. The graph constrains the relational structure, while the random vectors provide a reproducible numerical encoding of that structure. A fixed vector seed records one particular assignment within this scheme, and alternative seeds provide other numerical realizations of the same specified lexical structure, as explained in Section~\ref{sec:graph-encoding}. Vector construction precedes the neural realization, and the simulation examines how the resulting representations can be retrieved and mapped between lexical levels through spiking population activity. Learning these representations and associations from experience is a direction for future development.
